% Options for packages loaded elsewhere
\PassOptionsToPackage{unicode}{hyperref}
\PassOptionsToPackage{hyphens}{url}
\PassOptionsToPackage{dvipsnames,svgnames,x11names}{xcolor}
\documentclass[
]{article}
\usepackage{xcolor}
\usepackage[margin = 0.5in]{geometry}
\usepackage{amsmath,amssymb}
\setcounter{secnumdepth}{5}
\usepackage{iftex}
\ifPDFTeX
  \usepackage[T1]{fontenc}
  \usepackage[utf8]{inputenc}
  \usepackage{textcomp} % provide euro and other symbols
\else % if luatex or xetex
  \usepackage{unicode-math} % this also loads fontspec
  \defaultfontfeatures{Scale=MatchLowercase}
  \defaultfontfeatures[\rmfamily]{Ligatures=TeX,Scale=1}
\fi
\usepackage{lmodern}
\ifPDFTeX\else
  % xetex/luatex font selection
\fi
% Use upquote if available, for straight quotes in verbatim environments
\IfFileExists{upquote.sty}{\usepackage{upquote}}{}
\IfFileExists{microtype.sty}{% use microtype if available
  \usepackage[]{microtype}
  \UseMicrotypeSet[protrusion]{basicmath} % disable protrusion for tt fonts
}{}
\makeatletter
\@ifundefined{KOMAClassName}{% if non-KOMA class
  \IfFileExists{parskip.sty}{%
    \usepackage{parskip}
  }{% else
    \setlength{\parindent}{0pt}
    \setlength{\parskip}{6pt plus 2pt minus 1pt}}
}{% if KOMA class
  \KOMAoptions{parskip=half}}
\makeatother
\usepackage{longtable,booktabs,array}
\newcounter{none} % for unnumbered tables
\usepackage{calc} % for calculating minipage widths
% Correct order of tables after \paragraph or \subparagraph
\usepackage{etoolbox}
\makeatletter
\patchcmd\longtable{\par}{\if@noskipsec\mbox{}\fi\par}{}{}
\makeatother
% Allow footnotes in longtable head/foot
\IfFileExists{footnotehyper.sty}{\usepackage{footnotehyper}}{\usepackage{footnote}}
\makesavenoteenv{longtable}
\usepackage{graphicx}
\makeatletter
\newsavebox\pandoc@box
\newcommand*\pandocbounded[1]{% scales image to fit in text height/width
  \sbox\pandoc@box{#1}%
  \Gscale@div\@tempa{\textheight}{\dimexpr\ht\pandoc@box+\dp\pandoc@box\relax}%
  \Gscale@div\@tempb{\linewidth}{\wd\pandoc@box}%
  \ifdim\@tempb\p@<\@tempa\p@\let\@tempa\@tempb\fi% select the smaller of both
  \ifdim\@tempa\p@<\p@\scalebox{\@tempa}{\usebox\pandoc@box}%
  \else\usebox{\pandoc@box}%
  \fi%
}
% Set default figure placement to htbp
\def\fps@figure{htbp}
\makeatother
\setlength{\emergencystretch}{3em} % prevent overfull lines
\providecommand{\tightlist}{%
  \setlength{\itemsep}{0pt}\setlength{\parskip}{0pt}}
\usepackage{booktabs}
\usepackage{longtable}
\usepackage{array}
\usepackage{multirow}
\usepackage{wrapfig}
\usepackage{float}
\usepackage{colortbl}
\usepackage{pdflscape}
\usepackage{tabu}
\usepackage{threeparttable}
\usepackage{threeparttablex}
\usepackage[normalem]{ulem}
\usepackage{makecell}
\usepackage{xcolor}
\usepackage{titling}
\usepackage{bookmark}
\IfFileExists{xurl.sty}{\usepackage{xurl}}{} % add URL line breaks if available
\urlstyle{same}
\hypersetup{
  colorlinks=true,
  linkcolor={Maroon},
  filecolor={Maroon},
  citecolor={Blue},
  urlcolor={blue},
  pdfcreator={LaTeX via pandoc}}

\author{}
\date{\vspace{-2.5em}}

\begin{document}

\begin{titlepage}
\centering
\vspace*{\fill}

{\fontsize{24pt}{28pt}\selectfont \textbf{Classifying Power Plant Fuel Types\\[0.3em] via Emissions Fingerprinting}}

\vspace{1.5em}

{\large \textit{A Supervised Gaussian Naive Bayes Analysis of 2007 EPA eGRID Data}}

\vspace{3em}

{\Large \textbf{Alexander Elamthuruthy}}

\vspace*{\fill}
\end{titlepage}

\newpage
\pagenumbering{arabic}
\setcounter{page}{1}

\section{Summary}\label{summary}

Identifying the generation source of power facilities is critical for
emissions monitoring, environmental compliance, and regional grid
modeling. This study evaluates whether a power plant's primary fossil
fuel type (Coal, Natural Gas, Oil, or Other Fossil) can be classified
strictly through its annual output of four core air emissions:
\(\text{CO}_2\), \(\text{SO}_2\), \(\text{NO}_x\), and \(\text{CH}_4\).

Using plant-level records from the EPA's Emissions \& Generation
Resource Integrated Database (eGRID2007), raw skewed emissions features
were normalized via a \(\log(1 + x)\) transformation to align with
Gaussian distributional assumptions. A Gaussian Naive Bayes model
trained on a 70/30 split achieved an overall test accuracy of 77.6\%.
The classifier demonstrated strong predictive capability in isolating
coal-fired generation due to pronounced \(\text{SO}_2\) profiles, while
highlighting the vulnerability of probabilistic classifiers to severe
real-world class imbalances when distinguishing minority fuel
categories.

\section{Description of the Data}\label{description-of-the-data}

This data set belongs to the Agriculture/Environment/Ecology/Climate
data area. Specifically, it provides data about the environmental
characteristics for most of the electric power generated in the United
states for the year 2007. The data includes emissions, emission rates,
generation, heat input, resource mix, and may other attributes.

The objective of this study is to determine the extent to which a power
plant's primary fuel source can be identified solely through its
emission fingerprint. This research seeks to quantify the predictive
power of \(NO_x\), \(SO_2\), \(CO_2\), and \(CH_4\) outputs in
distinguishing between the different fuel sources used at power plants.

Data for this analysis was sourced from the U.S. Environmental
Protection Agency's (EPA) Emissions \& Generation Resource Integrated
Database (eGRID). Specifically, observations were extracted from the
2007 plant-level dataset (PLNT07), which compiles annual environmental,
operational, and generation metrics across power facilities operating
within the United States.

For this study we will be using the variables

\begin{itemize}
\tightlist
\item
  \textbf{PLFUELCT}: Primary Fossil Fuel Category, the categorical
  target variable (Coal, Gas, Oil, or Other Fossil).
\item
  \textbf{PLCO2AN}: Annual \(CO_2\) Emissions: Plant Annual carbon
  dioxide emitted (tons).
\item
  \textbf{PLSO2AN}: Annual \(SO_2\) Emissions: Plant Annual sulfur
  dioxide emitted (tons).
\item
  \textbf{PLNOXAN}: Annual \(NO_x\) Emissions: Plant Annual nitrogen
  oxides emitted (tons).
\item
  \textbf{PLCH4AN}: Annual \(CH_4\) Emissions: Plant Annual methane
  emitted (tons).
\end{itemize}

\section{Analysis \& Discussion}\label{analysis-discussion}

To find the predictive power of the emissions a Naive Bayes model was
used for classification. The response variable was PLFUELCT, the primary
fossil fuel used at the plant that had a value of either Coal, Gas, Oil,
or Other Fossil Fuel.

The Predictors were:

\begin{itemize}
\tightlist
\item
  \textbf{PLCO2AN}: Annual \(CO_2\) Emissions
\item
  \textbf{PLSO2AN}: Annual \(SO_2\) Emissions
\item
  \textbf{PLNOXAN}: Annual \(NO_x\) Emissions
\item
  \textbf{PLCH4AN}: Annual \(CH_4\) Emissions
\end{itemize}

\begin{longtable}[]{@{}
  >{\centering\arraybackslash}p{(\linewidth - 8\tabcolsep) * \real{0.1392}}
  >{\centering\arraybackslash}p{(\linewidth - 8\tabcolsep) * \real{0.2152}}
  >{\centering\arraybackslash}p{(\linewidth - 8\tabcolsep) * \real{0.2152}}
  >{\centering\arraybackslash}p{(\linewidth - 8\tabcolsep) * \real{0.2152}}
  >{\centering\arraybackslash}p{(\linewidth - 8\tabcolsep) * \real{0.2152}}@{}}
\caption{Plant means by Fuel Category}\tabularnewline
\toprule\noalign{}
\begin{minipage}[b]{\linewidth}\centering
PLFUELCT
\end{minipage} & \begin{minipage}[b]{\linewidth}\centering
Avg CO2 (Tons)
\end{minipage} & \begin{minipage}[b]{\linewidth}\centering
Avg SO2 (Tons)
\end{minipage} & \begin{minipage}[b]{\linewidth}\centering
Avg NOx (Tons)
\end{minipage} & \begin{minipage}[b]{\linewidth}\centering
Avg CH4 (Tons)
\end{minipage} \\
\midrule\noalign{}
\endfirsthead
\toprule\noalign{}
\begin{minipage}[b]{\linewidth}\centering
PLFUELCT
\end{minipage} & \begin{minipage}[b]{\linewidth}\centering
Avg CO2 (Tons)
\end{minipage} & \begin{minipage}[b]{\linewidth}\centering
Avg SO2 (Tons)
\end{minipage} & \begin{minipage}[b]{\linewidth}\centering
Avg NOx (Tons)
\end{minipage} & \begin{minipage}[b]{\linewidth}\centering
Avg CH4 (Tons)
\end{minipage} \\
\midrule\noalign{}
\endhead
\bottomrule\noalign{}
\endlastfoot
COAL & 3782559 & 15851 & 5536 & 87768 \\
GAS & 292271 & 41.98 & 153.7 & 11933 \\
OFSL & 117629 & 112 & 219.5 & 9103 \\
OIL & 48807 & 240.9 & 94.67 & 3467 \\
\end{longtable}

As evidenced by the summary statistics in Table 1, the raw emission
means show a massive difference in scale. For example, the average
\(CO_{2}\) output for Coal plants is measured in the millions of tons,
while the corresponding \(SO_{2}\) and \(NO_{x}\) averages are much
smaller in comparison. This extreme variance creates a heavily
right-skewed distribution that violates the fundamental assumptions of a
Gaussian Naive Bayes classifier. To address this, we applied a log
transformation (\(log1p\)) to all emission features.

\begin{center}\includegraphics{./figures/unnamed-chunk-3-1} \end{center}

Looking at the violin plot we can see the distribution of data values
for all the emissions based on the fuel used at a plant. When looking at
the plots we can predict that the model will be very accurate at
identifying Coal due to the separation of coal from the other fuels in
the \(SO_2\) plot. We can also predict that the model may struggle with
predicting oil due to the similar shape is has in all four of the plots

After cleaning the data and performing the log transformation, the data
was split into a training set and a test set, the training set holding
70\% of the data with the test set holding the other 30\%.

After fitting a model with the training set, we used said model to get
predictions of what the fossil fuel type should be based off of the
emissions. The predicted values were then compared to the test set in a
confusion matrix.

\begin{center}\includegraphics[width=0.5\textwidth]{./figures/unnamed-chunk-5-1} \end{center}

The model achieved an overall classification accuracy of 77.65\%. While
this figure suggests a generally effective model, a deeper examination
of the confusion matrix reveals significant differences in performance
across different fuel categories.

The model demonstrated its highest reliability in identifying Coal,
strengthening the prediction that the difference in \(SO_2\) compared to
the other fuels would make it identifiable. However, the model struggled
significantly with the remaining categories, most notably Other Fossil
Fuel (OFSL), which yielded a 0\% success rate.

This failure is primarily attributed to Representational Bias within the
dataset. As shown in the matrix, the Gas category contains over 300
observations, whereas OFSL contains fewer than 15. In a Gaussian Naive
Bayes framework, this imbalance causes the model to have a very poorly
defined ``bell curve'' for rare classes, making it mathematically more
likely to default to a more frequent class like Gas or Coal when
emissions data is ambiguous.

Consequently, while the model serves as a tool for identifying
high-emission coal facilities, its lack of sensitivity toward fuel types
with similar emissions makes it an unreliable standalone predictor for
the broader fuel mix.

\newpage

\section{Citations}\label{citations}

U.S. Environmental Protection Agency (EPA). Emissions \& Generation
Resource Integrated Database (eGRID), 2007 Plant Data. Washington, DC:
U.S. EPA.

\end{document}
